\documentclass{article}
\title{Segunda entrega de Proyecto de Análisis y Diseño de Software}
\author{Rodrigo Díaz-Regañón, Daniel Martínez Fernández\\ y Matteo Artuñedo González}
\date{\today}


\begin{document}
\maketitle
\tableofcontents
\newpage
\section{Diagrama de clases}
\subsection{Explicación del paquete usuario}
Todos los usuarios tienen unos atributos comunes: dni, nombre de usuario y contraseña, que aparecen como atributos en la clase abstracta Usuario.\\

Los usuarios se distinguen en dos grandes grupos, los clientes y los usuarios de gestión, que engloban al gestor y los empleados.\\

Con respecto a la clase UsuarioGestion, esta presenta todos los métodos destinados a la gestión de la tienda que podrán ser realizados tanto por los empleados como por el gestor. \\
\begin{itemize}
  \item Aparecen todos los métodos relativos a los productos. Estos métodos, o bien crearán un producto nuevo mediante su constructor e invocarán el método de la aplicación que añadirá el producto, o solicitarán a la aplicación el producto que deseen modificar y cambiaran sus atributos mediante setters y getters.
  \item Del mismo modo, tendrán métodos relativos a las categorías, con los que podrán añadirlas, eliminarlas y modificarlas.
  \item Lo mismo se sigue para los descuentos.
\end{itemize}
El gestor, en concreto, tendrá métodos con los cuales podrá:
\begin{itemize}
  \item Gestionar los empleados de la tienda, cambiando sus permisos y creándolos/eliminándolos.
  \item Configurar el sistema de recomendaciones, determinando la relevancia que se da a las novedades, las puntuaciones y los intereses de los usuarios, y configurando el número de productos que se recomendarán.
  \item Generar estadísticas de distinto tipo.
\end{itemize}
Con respecto al empleado, todos sus métodos verifican que el empleado tiene los métodos necesarios para ejecutarlos y, acto seguido, llaman a los métodos propios de las solicitudes para llevar a cabo las acciones. Propiamente en los métodos del empleado solo tiene lugar la comprobación de los permisos.

El cliente registrado, por su parte, cuenta con varios métodos agrupados por sus funcionalidades:
\begin{itemize}
  \item Para ordenar los clientes por gasto total en las estadística, se facilita un método que obtiene el gasto total del cliente a partir de sus pedidos y las validaciones de sus productos de segunda mano.
  \item El cliente podrá añadir y eliminar productos de su carrito y su cartera.
  \item Podrá gestionar ofertas. Las ofertas se realizan mediante el método realizarOferta, el cual llama al constructor de Oferta y la crea. Acto seguido, se introduce la oferta en los arrays correspondientes del usuario ofertante y el destinatario de la oferta. Aparte de la creación de ofertas nuevas, podrán aceptarlas o rechazarlas. Las ofertas rechazadas no se guardarán, únicamente aquellas que sean aceptadas y se transformen en una solicitud de intercambio.
  \item Podrá escribir reseñas de productos y eliminarlas. Para escribir una reseña, el producto deberá hallarse en su historial de productos comprados.
  \item A través de los métodos borrarNotificacion y añadirNotificacion la aplicación podrá insertar en el array de cada cliente una notificación determinada.
  \item Podrá pagar pedidos y validaciones. En ambos métodos, se realizará una llamada al método procesarPago de SistemaPago. Si el pago es exitoso, se devolverá un objeto de tipo Pago, que se asociará al producto de segunda mano o pedido según corresponda. Así quedará registrado el pago de los productos y los clientes.
\end{itemize}




\subsection{Explicación del paquete usuario}

El paquete de la clase aplicación es el paquete central de la aplicación. La clase Aplicación es la clase que contiene toda la información y, en caso de que se borrase, se borrarían todos los datos (asociación de contención fuerte) de la aplicación.\\

Los métodos de la clase aplicación se pueden agrupar de la siguiente forma:
\begin{itemize}
  \item Gestión de la cuenta: la aplicación, aparte del listado de todos los usuarios registrados, tiene una asociación de cardinalidad 0..1 a un usuario concreto. Esta referencia es nula cuando se inicia la aplicación, pero una vez se inicia sesión mediante el método correspondiente, pasa a apuntar al usuario que ha realizado el registro. El método para cerrar sesión deshace esta acción.
  \item Métodos que añaden y eliminan empleados del array de usuarios de la aplicación. Estos métodos son invocados por el gestor desde sus métodos propios.
  \item Métodos para la clientes relativos a productos a la venta. Estos incluye un método de búsqueda con una palabra clave (prompt), para el cual se accede
  \item Métodos para la creación de notificaciones y su envío a un usuario. El método para esto último será invocado desde métodos como aceptarOferta o validarProducto para informar a los usuarios de que estos eventos han tenido lugar.
  \item Métodos de pagos, que sirven como interfaz entre el sistema de pagos y los métodos relativos a los pagos que se encuentran en las clases
\end{itemize}
La aplicación tiene como módulo un gestor de solicitudes, que se encarga de añadir solicitudes a la aplicación para tener registro de todas ellas.
Existe también como módulo un sistema de pago, que se encarga de procesar los pagos con tarjeta.
Asimismo, existe un sistema de estadísticas que permite generar estadísticas de diferente tipo para unos períodos dados. \\
Asociada a la aplicación está también la clase ConfiguracionRecomendacion, que guarda la configuración que el gestor desee determinar de cara a recomendar productos a los empleados.
Por último, se encuentra el Catálogo. Esta clase es la que guarda todos los productos, las categorías y los descuentos de la tienda, así como los filtros aplicados para la búsqueda. Contiene métodos para:
\begin{itemize}
  \item Gestionar los productos de la tienda y modificarlos.
  \item Gestionar las categorías de la tienda y modificarlas
  \item Aplicar descuentos
\end{itemize}

\subsection{Explicación del paquete producto}
EL paquete de productos agrupa todos los artículos que se pueden gestionar en la aplicación, ya sean productos 
destinados a la venta oficial de la tienda o artículos de segunda mano aportados por los clientes para el intercambio.
Toda clase de producto comparte los atributos:
\begin{itemize}
  \item id: el identificador único que permite distinguir los artículos en el sistema.
  \item name: el nombre del producto.
  \item description: un texto que explica y detalla las características del artículo.
  \item photo: el archivo de imagen que mostrará el aspecto del producto en la interfaz.
\end{itemize}

Existe una clase genérica de producto que únicamente contiene los atributos ya mencionados.
A partir de ella, los productos se distinguen en los destinados a la venta y aquellos orientados al intercambio.

\subsubsection{Productos de venta}
Los productos que vende directamente la tienda (LineaProductoVenta) incluyen su stock disponible, su precio y la cantidad de unidades vendidas. Estos podrán:
\begin{itemize}
  \item Recibir reseñas: los clientes registrados podrán dejar opiniones (Reseña) que guardarán un texto descriptivo, una puntuación y la fecha de publicación.
  \item Formar packs: un producto de venta podrá estar compuesto a su vez por una agrupación de varios productos.
  \item Incluirse en compras: cuando un cliente los añade a su carrito se convierte en ProductoVenta, registrando las unidades exactas que se van a comprar para enlazarse con una solicitud de pedido.
  \item Tener categorias: los productos pueden estar clasificados en categorias, y cada categoria puede tener muchos productos.
\end{itemize}

A su vez, los artículos del catálogo se dividen en tres tipos concretos:
\begin{itemize}
  \item Juegos de mesa: registran el número de jugadores, las edades mínima y máxima recomendadas, y su formato exacto (Cartas, Dados o Miniaturas).
  \item Cómics: guardan el número de páginas, su autor, la editorial y su año de publicación.
  \item Figuras: detallan su marca, el material de fabricación y sus tres dimensiones físicas (X, Y, Z).
\end{itemize}

  \subsubsection{Productos de segunda mano}
Los artículos que suben los clientes para intercambiar (ProductoSegundaMano) cuentan con un atributo booleano (validado) que indica si ya han recibido la revisión de la tienda. 

Al ser revisados por un empleado, los productos de segunda mano ya tienen una validacion. Se les adjuntan los datos de dicha revisión (DatosValidacion), que contienen un precio estimado y 
clasifican el estado de conservación del artículo (Perfecto, Muy bueno, Uso ligero, Uso evidente, Muy usado o Dañado).
Estos productos podrán ser ofertados o solicitados por los clientes mediante ofertas.

\subsection{Explicación del paquete descuento}
EL paquete de descuento agrupa todos los descuentos que pueden aplicarse en la aplicación.
Todo descuento comparte los atributos:
\begin{itemize}
  \item fechaInicio: momento en el que el descuento entra en vigor.
  \item fechaFin: momento en el que el descuento deja de estar disponible.
\end{itemize}

Existe un tipo genérico de descuento que contiene los atributos ya mencionados y que se asocia a los elementos rebajados, pudiendo aplicarse a una lista de productos concretos (LineaProductoVenta) o a una categoría (Categoria).
Los descuentos se distinguen entre aquellos que se aplican directamente y los que requieren alcanzar un gasto mínimo.

\subsubsection{Directos}
Los descuentos directos podrán ser de los siguientes tipos:
\begin{itemize}
  \item Descuentos de cantidad: en caso de ofertas por volumen (tipo 3x2), el descuento (Cantidad) guardará la cantidad de productos requerida (numeroComprados) y la cantidad final que se lleva el cliente (numeroRecibidos).
  \item Descuentos de precio: el descuento (Precio) aplicará una reducción según el porcentaje indicado en su atributo (porcentaje).
\end{itemize}

\subsubsection{Umbral de gasto}
Los descuentos por umbral de gasto (UmbralGasto) tendrán todos un importe mínimo (umbral) que el cliente debe superar en su compra, y se corresponderán con:
\begin{itemize}
  \item Una rebaja: al alcanzar el umbral, se aplicará un porcentaje de descuento extra en el total de la compra (Rebaja), guardado en su atributo (porcentajeRebaja).
  \item Un regalo: al alcanzar el umbral, se otorgará al cliente un producto gratuito (Regalo), desde el cual se podrá acceder al producto concreto que se regala (ProductoVenta).
\end{itemize}

\subsection{Explicación del paquete categoría}
EL paquete de categoría agrupa las diferentes clasificaciones que permiten organizar el catálogo de la aplicación.
Toda categoría comparte el atributo:
\begin{itemize}
  \item nombre: el nombre que permite distinguir y etiquetar las categorías entre sí.
\end{itemize}

Existe una clase central de categoría que únicamente contiene el atributo ya mencionado.\\

Las categorías se relacionan con otros paquetes del sistema, distinguiéndose entre los productos que agrupan y los descuentos que se les aplican.

\subsubsection{Productos}
Las categorías podrán organizar los siguientes elementos:
Líneas de producto: para mantener el catálogo estructurado, una categoría podrá contener múltiples productos, y a su vez, 
desde un producto (LineaProductoVenta) se podrá acceder a todas las categorías a las que pertenece.

\subsubsection{Descuentos}
Las categorías podrán tener descuentos asociados, que se corresponderán con:
\begin{itemize}
  \item Un descuento: una categoría podrá tener un descuento activo (o ninguno, de ahí la relación 0..1), mientras que un mismo descuento genérico (Descuento) podrá aplicarse simultáneamente a múltiples categorías. 
La información relativa a estos descuentos se explicó con más detalle en el paquete de los descuentos.
\end{itemize}

\subsection{Explicación del paquete notificación}
EL paquete de notificación agrupa todas las notificaciones que pueden recibir los usuarios de la aplicación.
Toda notificación comparte los atributos:
\begin{itemize}
  \item id: el indentificador que permite distinguir las notificaciones entre sí. Se asigna según el orden de creación, de modo que la primera 
notificación en crearse tiene el id 1.
\item mensaje: el mensaje que le aparecerá al usuario que reciba la notificación, el cual explicará el contenido de esta.
\item horaEnvio: se guardará el instante preciso en el que se envió la notificación, para poder acceder a él en caso de que surgiese algún incidente con un cliente o empleado (por ejemplo, que el cliente reclame que la notificación no le llegó a tiempo para un intercambio).
\end{itemize}

Existe un tipo genérico de notificación que únicamente contiene los atributos ya mencionados.Este tipo de notificaciones se podrán
emplear para enviar mensajes genéricos, como informar a todos los clientes y empleados de que la tienda cierra por vacaciones.

Las notificaciones concretas se distinguen entre aquellas que pueden recibir los clientes registrados y las que pueden recibir los empleados. 


\subsubsection{Cliente}
Los clientes podrán recibir las siguientes notificaciones:
\begin{itemize}
  \item Notificaciones de productos: en caso de que haya productos novedosos o rebajados, el cliente podrá recibir una notificación (NotificacionProducto) con la que podrá acceder a todos estos productos (lista de objetos de tipo LineaProductoVenta).
\item Notificaciones de oferta: en caso de que el cliente reciba una oferta o se le acepte/rechace una oferta que ya ha enviado, recibirá una notificación (NotificacionOferta) desde la cual podrá acceder a la oferta concreta.
\item Notificaciones de validación: si un producto de un cliente es validado por un empleado, recibirá una notificación (NotificacionValidacion)desde la cual podrá acceder al producto, que ahora tiene información sobre su validación.
\item Notificaciones de pedido: cuando se cambie el estado de un pedido, un cliente recibirá una notificación (NotificacionPedido) 
desde la cual podrá acceder al pedido que ha realizado.
\item Notificaciones de intercambio: una vez se haya confirmado un intercambio, el cliente recibirá una notificación con los datos
de dicho intercambio.
\end{itemize}

\subsubsection{Empleado}
Las notificaciones de los empleados tendrán todas una solicitud asociada, que  se correspondrá con:
\begin{itemize}
  \item Un pedido que el empleado puede gestionar, cambiando sus estados según se vaya preparando y se entregue.
  \item Una validación que el empleado puede realizar, introduciendo los datos de la validación.
  \item Un intercambio que el empleado puede aprobar presencialmente.
\end{itemize}
La información relativa a las solicitudes se explicará con más detalle en el paquete de las solicitudes.

\subsection{Paquete solicitud}

EL paquete de solicitud agrupa todas las peticiones y trámites que los usuarios realizan en el sistema.

Existe una clase abstracta de solicitud (Solicitud) de la que heredan todos los tipos de gestiones.
Las solicitudes se distinguen en tres tipos principales:

\subsubsection{Pedidos}
Los pedidos (SolicitudPedido) representan las compras en bloque (en lugar de producto por producto) que se realizan en la tienda y están vinculados al cliente que los realiza (ClienteRegistrado). Estos podrán:
\begin{itemize}
  \item Incluir productos: agrupan los productos de venta que el cliente va a adquirir (ProductoVenta).
  \item Registrar un pago: pueden tener asociado un comprobante económico (Pago) que guarda su fecha (fechaPago) y el coste total (importe). De este modo, tanto el empleado como el cliente pueden verificar que el pago 
ha sido realizado y ha quedado asociado al pedido correctamente.
\item Seguimiento de estado: su progreso se monitoriza mediante un indicador (EstadoPedido) que refleja en qué punto está la compra (Pendiente de pago, Pagado, Listo para recogida o Recogido). Será el empleado el encargado de modificar estos datos según se progrese con el pedido en la tienda.
\end{itemize}
El pedido se crea una vez el cliente decide pedir los productos de su carrito. El objeto se asocia a los pedidos del cliente y se guarda también en el array de pedidos de la aplicación. Será este array global el que podrán solicitar los empleados con permisos para buscar en él los pedidos cuyo estado quieran cambiar e invocar los métodos pertinentes. El cliente también podrá cancelar el pedido si lo desease y pagarlo.

\subsubsection{Validaciones}
Las solicitudes de validación (SolicitudValidacion) se corresponden a peticiones de los clientes para que la tienda valore sus productos de segunda mano. Estas solicitudes tendrán asociadas:
\begin{itemize}
  \item Un producto a validar: apuntan directamente al artículo específico (ProductoSegundaMano) que el empleado debe evaluar para certificar su estado.
  \item Un pago de revisión: al igual que los pedidos, pueden tener un pago asociado (Pago), al que podrán acceder el cliente (a través del producto de segunda mano existe navegación a la solicitud, que tiene asociado el pago) y el empleado (que podrá ver el listado de todas las solicitudes a través de la clase Aplicación).
  \item La solicitud de validación se crea en el momento en el que el usuario sube el producto de segunda mano a la tienda. Se asocia la solicitud al producto y se guarda también en la lista de todas las solicitudes de validación que posee la aplicación.
\end{itemize}
Así, cuando un empleado quiera realizar una validación, podrá solicitar el array de las validaciones a la aplicación, buscar en él la solicitud que desee e invocar sobre ella los métodos pertinentes.

\subsubsection{Intercambios}
Las solicitudes de intercambio (SolicitudIntercambio) gestionan el intercambio en tienda de productos de segunda mano.
Todo intercambio está vinculado a una propuesta previa (Oferta), la cual registra qué cliente es el ofertante, quién es el destinatario, junto a los códigos que deberán ser presentados en tienda, y qué productos de segunda mano se ofrecen y se solicitan. En el momento en el que se acepta una oferta, se crea la solicitud de intercambio y quedan vinculadas. En la solicitud se guardarán: los detalles del intercambio, que contiene la información del evento final (DetallesIntercambio), concretando la fecha exacta (fechaIntercambio) y el lugar de encuentro (lugarIntercambio) en caso de que hubiese varias tiendas físicas.
Los clientes también reciben estos detalles a través de una notificación (NotificacionIntercambio).
Cabe mencionar que los clientes no podrán ejecutar los getters de los códigos, solo lo podrá hacer el empleado, garantizando así la seguridad del intercambio.


\subsection{Paquete filtro}
El paquete filtro, nace para simplificar el filtrado en las búsquedas tanto de productos de la tienda y productos de intercambio, como productos que se quieran modificar, no contiene métodos especiales más allá de getters y setters, ya quesu función es ir almacenando los filtros que el usuario ha seleccionado durante una sesión y estos se resetearan al terminar la misma.
Todo filtro comparte un atributo: ordenAscendente, de tipo boolean que determinará si los productos se verán en orden ascendente o descendente, siguiendo un atributo que se especificará en el método de búsqueda.

\subsubsection{Cliente registrado}
Podemos hacer la primera distinción entre productos de venta y de intercambios, estos últimos contienen el valor mínimo y valor máximo que 
determinaran el rango de valoración por el que el cliente quiere buscar, así como el estado de los productos que le interesa al cliente, pudiendo
elegir cuales quiera de ellos simultáneaente.

\subsubsection{Usuario gestión}
Por otro lado, los productos de venta contienen las categorías que se desean buscar, así como el tipo de producto, pudiendo elegirse varios tipos
o categorías simultáneamente, estos son los filtros de los que disponen los usuarios que puedan gestionar productos. 

\subsubsection{Cliente}
Existe otro tipo de filtros que hereda de los filtros de venta, estos son los empleados por los clientes a la hora de filtrar sus búsquedas. Además de todo lo anterior, contienen el intervalo de precio deseado (con un precio mínimo y un precio máximo), así como un rango de puntuación (con una puntuación mínima y una máxima), así como un atributo de tipo boolean (ordenarPorPrecio), que en caso de ser TRUE se ordenará la búsqueda por precio, de lo contrario se ordenará por puntuación. Por último contendrá la selección del tipo de descuento en los productos, se podrán seleccionar varios tipos de descuentos simultáneamente. 

\section{Diagramas de secuencia}
\subsection{PagarPedido}
Este diagrama describe un proceso crítico de negocio que involucra transacciones financieras y gestión de estados de pedidos.
\subsubsection{Fase de Cálculo y Procesamiento}
\begin{itemize}
\item Cálculo de Costes: La Aplicacion interactúa con la SolicitudPedido (p) para obtener el coste total. Esta última recorre sus líneas de producto para sumar los precios individuales, ilustrando una delegación de responsabilidades correcta.
\item Pasarela de Pago: Se invoca al SistemaPago mediante el método procesarPago(...) pasando los datos de la tarjeta y el coste calculado. El resultado (validez) determina el flujo siguiente.
\end{itemize}
\subsubsection{Gestión de Excepciones y Notificaciones}
\begin{itemize}
\item Validación Negativa (opt): Si el pago falla (validez == FALSE), se utiliza un fragmento opt para manejar el error (aunque el diagrama se centra más en el flujo de éxito posterior).
\item Confirmación y Notificación: Si el pago es exitoso:
  \begin{itemize}
    \item Se actualiza el estado del pedido a PAGADO.
    \item Se genera una notificación para el empleado.
    \item Bucle de Envío (loop y break): Se intenta enviar la notificación hasta que noti == TRUE, momento en el cual el fragmento break interrumpe el bucle de reintentos. Esto se traduce como intentar asignar un pedido a un empleado hasta que se encuentra a uno libre.
  \end{itemize}
\end{itemize}
\subsubsection{Post-procesamiento}
\begin{itemize}
  \item Actualización del Interés: Al igual que en el diagrama de búsqueda, el proceso finaliza actualizando el perfil de interés del cliente basándose en los productos efectivamente comprados. Esto asegura que el sistema "aprenda" de las transacciones finalizadas.
\end{itemize}
\subsection{BuscarProductosNuevos}
Este diagrama modela la lógica de una búsqueda filtrada, integrando controles de seguridad y personalización del interés del usuario.
\subsubsection{Flujo de Control y Seguridad}
\begin{itemize}
  \item Validación de Usuario (break): El diagrama inicia verificando la identidad del usuarioActual. Se utiliza un fragmento break con la guarda [clase != null AND clase != ClienteRegistrado]. Si el usuario no es un cliente registrado, el proceso se interrumpe inmediatamente devolviendo FALSE, lo que garantiza que solo usuarios autorizados accedan a esta funcionalidad específica.
\item Obtención del Filtro: La Aplicacion solicita al Catalogo los productos, pero primero debe obtener un objeto FiltroVenta para aplicar los criterios de búsqueda.
\end{itemize}
\subsubsection{Lógica de Filtrado Detallada}
El segundo bloque del diagrama profundiza en el método pasaFiltro de la clase LineaProductoVenta:
\begin{itemize}
  \item Recuperación de Parámetros: El producto consulta al objeto filtro para obtener los umbrales de puntuación, precio y tipos de descuento permitidos.
  \item Evaluación de Condiciones (alt): Se utiliza un fragmento alt con una guarda compleja que verifica simultáneamente:
  \begin{itemize}
    \item Rangos de precio y puntuación media.
    \item Contenencia del prompt en el nombre del producto.
    \item Validez del tipo de descuento aplicado.
  \end{itemize}
\item Resultado: Si cumple, el producto se añade a la colección de resultados; de lo contrario, se devuelve FALSE.
\end{itemize}
  \subsubsection{Post-procesamiento}
\begin{itemize}
  \item Actualización de Interés: Una vez obtenidos los productos, el sistema recupera el objeto Interes del usuario y llama a actualizarInteresBusquedaCompra(productos). Esto demuestra un diseño orientado a la personalización basado en el comportamiento de búsqueda.
\end{itemize}

\section{Diagramas de estado}
\subsection{Intercambio}
Este diagrama representa el ciclo de vida de un intercambio, mostrando su evolución desde que se propone hasta que se completa o se anula. 
Todo comienza cuando un cliente realiza una oferta, quedando el proceso en estado de solicitud. 
En esta fase inicial, si la oferta es rechazada, el intercambio se cancela directamente y termina su ciclo. 
Si, por el contrario, la contraparte acepta la oferta, el proceso avanza al estado de intercambio aceptado. 
Finalmente, se realiza una aprobación por parte del empleado; si esta validación falla, el intercambio también queda cancelado, pero si se aprueba con éxito, el ciclo concluye de forma quedando el intercambio oficialmente registrado.

\subsection{Pedido}
Este diagrama representa todos los estados por los que puede pasar un pedido, mostrando cómo avanza desde su creación hasta su entrega al cliente o su cancelación.
Todo comienza cuando se realiza el pedido, que queda a la espera de ser pagado; si el intento de pago falla, 
el pedido se mantiene en ese mismo punto, aunque ahora también existe la opción de cancelar el pedido definitivamente en esta fase, 
terminando ahí su ciclo. Si continúa y el sistema confirma el cobro, el pedido avanza al estado de "pagado". 
A partir de ahí, el empleado irá actualizando el estado del pedido (cuando esté listo para recogerse y finalmente al ser recogido por el cliente).
\end{document}
